\documentclass[11pt]{article}
\usepackage{mathblatt}
% Präambel des Setzers (werkzeuge/setzer-praeambel.tex, wortgleich); dazu \gteil (Muster T2U)
\makeatletter
\setlength{\mbpfbe}{0mm}% keine Punkte (pruefung.md § 5: nur beim Prüfstein)
\newlength{\gEin}\setlength{\gEin}{0pt}
\let\gSchrift\relax
\renewcommand{\pfkreuzzeile}[1]{\par\noindent\hangindent3.2em\hangafter1\hspace*{1.6em}$\square$\ #1\par}
\newcommand{\gkreuz}[1]{$\square$\ #1\hspace{2em}}
\newcommand{\gkreuzzeile}[2]{\par\noindent\hangindent2.2em\hangafter1\mbox{#1\,$\square$}\ #2\par}% Buchstabe, Option: eine je Zeile (neben Skizze)
% links, was man ansieht, rechts, was man tut; Spaltengrenze überall an derselben Stelle (gemeinsam.md „Satz“)
\newsavebox{\gbox}\newlength{\gLinks}
% v0.10 (Lehrer 13.10.): keine starre Halbierung – die linke Spalte ist so breit wie die Figur (+ 6 mm), die
% Antwort steht direkt daneben; ist die Figur breiter als 0,62 Zeile, steht die Antwort darunter
\newcommand{\gzwei}[2]{\par\smallskip\noindent\sbox{\gbox}{#1}%
  \setlength{\gLinks}{\dimexpr\wd\gbox+6mm\relax}%
  \ifdim\gLinks>0.62\linewidth
    \usebox{\gbox}\par\smallskip\noindent\begin{minipage}{\linewidth}#2\end{minipage}%
  \else
    \begin{minipage}[t]{\gLinks}\vspace{0pt}\usebox{\gbox}\end{minipage}%
    \begin{minipage}[t]{\dimexpr\linewidth-\gLinks\relax}\vspace{0pt}\raggedright #2\end{minipage}%
  \fi\par}
% v0.9 (12.10., pruefung.md § 5 „Marke mit B, Bedingungen rechts“): drittes Argument von pfaufg = Bedingungen
% der Prüfung (v0.10: \pfohnetr = TR durchgestrichen, „FOR“), klein grau im rechten Rand auf der Grundlinie der Nummer, Zeilen mit \\
\newcommand{\gbedrand}[2]{\ifblank{#2}{}{\rlap{\hspace*{\dimexpr#1+2mm\relax}%
  \parbox[t]{15mm}{\raggedright\scriptsize\leavevmode\color{mbgrau}\strut#2}}}}
% Teil eines Bündels (a, b …): Marke links im Rand, Bedingungen rechts im Rand, auf der Zeile des Teils
\newcommand{\gteilrand}[2]{\llap{\mbpfrandmarke{#1}\hspace*{\mbpfnr}}\gbedrand{\linewidth}{#2}}
% Teil eines Bündels: Buchstabe in eigener Spalte, Inhalt eingerückt (\gEin = Einzug, damit die Spaltengrenze von
% \gzwei an derselben Stelle bleibt)
\newenvironment{gbteil}[3]{\par\smallskip\noindent\gteilrand{#2}{#3}\makebox[\mbpfteil][l]{\textbf{#1}}%
  \begin{minipage}[t]{\dimexpr\linewidth-\mbpfteil\relax}\setlength{\gEin}{\mbpfteil}\raggedright\strut\ignorespaces}%
 {\end{minipage}\par}
% Grafik höchstens so breit wie die Spalte (Bündel nebeneinander, Koordinatensystem ohne Karo)
\newsavebox{\gpbox}
\newcommand{\gpassend}[1]{\sbox{\gpbox}{#1}\ifdim\wd\gpbox>\linewidth\typeout{GPASSEND: verkleinert \the\wd\gpbox\space auf \the\linewidth}\resizebox{\linewidth}{!}{\usebox{\gpbox}}\else\usebox{\gpbox}\fi}
% Skizzenfeld (v0.10, Lehrer 13.10.): Karo ab dem linken Rand über die Zeile, „Skizze“ klein grau oben links im Feld
% (wie \pfskizzenfeld der Vorlage); der rechte Rand bleibt den Bedingungen
\newcommand{\gskizze}[1]{\par\smallskip\noindent\begin{tikzpicture}
  \clip (0,0) rectangle ({\linewidth-1mm},{#1*10mm});
  \draw[black!16,line width=0.3pt,step=5mm] (0,0) grid ({\linewidth},{#1*10mm});
  \node[anchor=north west,inner sep=2pt,font=\scriptsize,text=mbgrau,fill=white] at (0,{#1*10mm}) {Skizze};
  \end{tikzpicture}\par}
\RenewDocumentEnvironment{pfaufg}{m m m}{%
  \begin{lrbox}{\mbaufgabebox}\begin{minipage}[t]{\linewidth}%
  \gSchrift\noindent\hspace*{\gEin}\mbpfrandmarke{#1}\gbedrand{\linewidth-\gEin}{#3}%
  \begin{minipage}[t]{\mbpfnr}\strut\textbf{#2}\end{minipage}%
  \begin{minipage}[t]{\dimexpr\linewidth-\gEin-\mbpfnr\relax}\raggedright\strut\ignorespaces}%
 {\end{minipage}%
  \end{minipage}\end{lrbox}%
  \setlength{\mbaufgabehoehe}{\dimexpr\ht\mbaufgabebox+\dp\mbaufgabebox\relax}%
  \ifdefempty{\gkopfnach}{\par\addvspace{5pt}\Needspace*{\mbaufgabehoehe}}{%
    % Stufenkopf, der mit dieser (ersten) Aufgabe zusammenbleibt: Platz für beide verlangen, dann den Kopf setzen
    % (12.10., Quadratische: ein Kopf stand allein unten, die Aufgabe mit Koordinatensystem auf der nächsten Seite)
    \par\Needspace*{\dimexpr\mbaufgabehoehe+40pt\relax}\gkopfnach\gdef\gkopfnach{}\par\addvspace{5pt}}%
  \noindent\usebox{\mbaufgabebox}\par\addvspace{5pt}}
% Stufenkopf vormerken: \gkopfvor{<Kopf>} – die nächste pfaufg setzt ihn, sobald beide zusammen Platz haben
\def\gkopfnach{}
\newcommand{\gkopfvor}[1]{\gdef\gkopfnach{#1}}
% Teilaufgabe mit Skizze links in gzwei: Buchstabe in eigener Spalte, darunter Text und Skizze
\newcommand{\gteil}[2]{\makebox[\mbpfteil][l]{\textbf{#1}}\parbox[t]{\dimexpr0.5\textwidth-\gEin-\mbpfnr-5mm-\mbpfteil\relax}{\raggedright #2}}
% Stufenkopf der Vorlage (mathblatt.sty 2026-10-12: Linie, Abstand, 13 pt, \pfjahre) ohne ihren Vorhalt von
% 14 Zeilen: der Setzer fordert vor dem Kopf 0,16 Seitenhöhe (Kopf + erste Aufgabe), mehr nicht
\setlength{\mbpfstufevorhalt}{0pt}
\makeatother
\newcommand{\karo}[1]{\par\smallskip\noindent\begin{tikzpicture}
  \clip (0,0) rectangle ({\linewidth-1mm},{#1*10mm});
  \draw[black!16,line width=0.3pt,step=5mm] (0,0) grid ({\linewidth},{#1*10mm});
  \end{tikzpicture}\par}
\newcommand{\kk}[1]{$\square$\,#1\hspace{0.9em}}
\newcommand{\linie}[1]{\,{\color{black!45}\rule[-1pt]{#1}{0.4pt}}\,}
\newcommand{\pkt}{\textperiodcentered{}}
% v0.14 (14.10.): Stufenkopf mit fünf Punkten (Option = Jahre für die Bestellkarte) oder „zuletzt 19“
\NewDocumentCommand{\gpunkte}{o m}{#2}
\newcommand{\gzuletzt}[1]{zuletzt #1}
% v0.7 (12.10.): Schülerfrage unter dem Stufenkopf; Ankreuzoption als Bild (Buchstabe und Kästchen darunter)
\newcommand{\gfrage}[1]{\par\nobreak\noindent\hspace*{\mbpfnr}{\small\itshape #1}\par\nobreak\addvspace{2pt}}
\newlength{\gbildbreite}\setlength{\gbildbreite}{0.32\linewidth}% Setzer setzt sie je Reihe (vier Bilder: 0,24)
\newcommand{\gbildoption}[3]{\begin{minipage}[t][#3][b]{\gbildbreite}\centering #2\par\smallskip\mbox{#1\,$\square$}\end{minipage}}
% Lineare (12.10.): Bild mit Buchstabe ohne Kästchen (Diagramm zuordnen); Aussage mit „wahr / falsch“ zum Ankreuzen
\newcommand{\gbildname}[3]{\begin{minipage}[t][#3][b]{\gbildbreite}\centering #2\par\smallskip\mbox{\textbf{#1}}\end{minipage}}
\newcommand{\gwahrfalsch}[2]{\par\noindent\hangindent2.2em\hangafter1\makebox[2.2em][l]{(#1)}#2\hfill\mbox{wahr $\square$\quad falsch $\square$}\par}
\makeatletter\setlength{\mbpfbe}{0mm}\makeatother
\begin{document}

\pfheftstil
\pflaufkopf[Abitur GK \pkt{} mehr]{Erwartungswert im Spiel}{AS90}\pfseitevon
\pfheftkopf{Erwartungswert im Spiel}{mehr}
\pfstufenzeile{Erwartungswert deuten \pkt{} Unbekannte gesucht}
% 2018-be-gk-B3.1d-m1 (Schwere I, rang 5)
\begin{pfaufg}{GK~’18\\B3.1d}{1.}{}
Bei einem Spiel beträgt der Einsatz 2 Euro. Mit der Wahrscheinlichkeit 0,3 werden 5 Euro ausgezahlt, sonst ist der Einsatz verloren. Zeige, dass das Spiel für den Anbieter auf lange Sicht gewinnbringend ist.
\karo{3}
\fusshilfe{\mbox{1:~Anbieter gewinnt im Mittel 0,50 Euro je Spiel}}
\end{pfaufg}
% 2018-iqb-ga-A12b-m1 (Schwere I, rang 6)
\begin{pfaufg}{IQB~’18\\A12b}{2.}{}
Die Auszahlung $X$ in Euro hat die Werte 1, 3 und 9 mit den Wahrscheinlichkeiten $\tfrac{1}{2}$, $\tfrac{1}{3}$ und $\tfrac{1}{6}$. Der Einsatz beträgt 3 Euro. Berechne den Erwartungswert von $X$ und deute ihn unter Berücksichtigung des Einsatzes.
\karo{3}
\fusshilfe{\mbox{2:~auf lange Sicht gleichen sich Einsatz und Auszahlung aus (faires Spiel)}}
\end{pfaufg}
% 2021-be-gk-A1.6b-m1 (Schwere II, rang 7)
\begin{pfaufg}{GK~’21\\A1.6b}{3.}{}
Ein roter und ein blauer Tetraeder (Zahlen 1 bis 4) werden geworfen. $X = 0$, wenn die rote Zahl kleiner als die blaue ist, $X = 1$ bei gleichen Zahlen und $X = 2$, wenn die rote Zahl größer ist. Ermittle den Erwartungswert von $X$.
\karo{3}
\fusshilfe{\mbox{3:~1}}
\end{pfaufg}
% 2024-iqb-ga-A12b-m1 (Schwere II, rang 8)
\begin{pfaufg}{IQB~’24\\A12b}{4.}{}
Bei einem Spiel zahlt man 3 Euro Einsatz und dreht ein Glücksrad zweimal; der grüne Sektor kommt mit 20 \%. Für jedes Grün gibt es 5 Euro. Die Wahrscheinlichkeit für genau einmal Grün ist 0,32. Begründe, dass der Veranstalter auf lange Sicht mehr einnimmt als auszahlt.
\karo{4}
\fusshilfe{\mbox{4:~Veranstalter nimmt mehr ein}}
\end{pfaufg}
% 2023-iqb-ga-A11a-m1 (Schwere I, rang 9)
\begin{pfaufg}{IQB~’23\\A11a}{5.}{}
Die Zufallsgröße $X$ nimmt die Werte 1, 2, 5 und $x_4$ mit den Wahrscheinlichkeiten $\tfrac{1}{2}$, $\tfrac{1}{4}$, $\tfrac{1}{8}$ und $\tfrac{1}{8}$ an; ihr Erwartungswert ist 3. Berechne $x_4$.
\karo{4}
\fusshilfe{\mbox{5:~11}}
\end{pfaufg}
% 2021-be-gk-B4i-m1 (Schwere II, rang 10)
\begin{pfaufg}{GK~’21\\B4i}{6.}{}
Beim täglichen Start eines Spiels gibt es 10 Bonuspunkte mit 60 \%, 20 Bonuspunkte mit 30 \% und 50 Bonuspunkte mit 10 \%. Die Wahrscheinlichkeiten für 10 und 20 Punkte werden so geändert, dass in 200 Tagen im Mittel 3 200 Bonuspunkte anfallen; 50 Punkte bleiben bei 10 \%. Ermittle die beiden neuen Wahrscheinlichkeiten.
\karo{5}
\fusshilfe{\mbox{6:~0,2}}
\end{pfaufg}
% 2025-iqb-ga-A22b-m1 (Schwere III, rang 11)
\begin{pfaufg}{IQB~’25\\A22b}{7.}{}
In einem Behälter liegen eine schwarze und $w$ weiße Kugeln ($w \ge 2$). Man zieht zweimal ohne Zurücklegen; Einsatz 3 Euro. Schwarz im ersten Zug bringt 9 Euro, schwarz im zweiten Zug 6 Euro, sonst nichts. Auf lange Sicht gleichen sich Einsätze und Auszahlungen aus. Ermittle $w$.
\karo{5}
\fusshilfe{\mbox{7:~4}}
\end{pfaufg}
% 2018-be-gk-B3.1e-m1 (Schwere III, rang 12)
\begin{pfaufg}{GK~’18\\B3.1e}{8.}{}
In einem Topf liegen 4 weiße und 3 schwarze Kugeln. Es kommen 11 weiße und $x$ schwarze hinzu. Zwei Kugeln werden mit einem Griff gezogen; man gewinnt, wenn keine schwarze dabei ist. Das Spiel ist fair, wenn die Gewinnwahrscheinlichkeit 0,5 beträgt. Zeige, dass $x = 2$ nicht zu 0,5 führt, und ermittle $x$ für ein faires Spiel.
\karo{7}
\fusshilfe{\mbox{8:~0,553 $\neq$ 0,5; $-$38 entfällt)}}
\end{pfaufg}
\par\vfill\noindent{\scriptsize\color{mbgrau}Originale: Abi GK 2018 \pkt{} B3.1d, IQB GK 2018 \pkt{} A12b, Abi GK 2021 \pkt{} A1.6b, IQB GK 2024 \pkt{} A12b, IQB GK 2023 \pkt{} A11a, Abi GK 2021 \pkt{} B4i, IQB GK 2025 \pkt{} A22b, Abi GK 2018 \pkt{} B3.1e\par}
\end{document}